% ============================================================
%   Appendix C  \textemdash\  Reproducibility Guide
% ============================================================

\chapter{Reproducibility Guide}\label{app:reproducibility}

This appendix lists the exact commands needed to reproduce every experiment reported in
Chapter~\ref{ch:experimentation}, together with the SQL queries used to extract aggregate
statistics from the per-run SQLite databases.

\section{Environment setup}\label{app:repro-setup}

\begin{enumerate}
  \item Install the CARLA 0.9.16 binary distribution. On Windows, launch the server with the
        command shown in Figure~\ref{fig:repro-launch}.
  \item Activate the project virtual environment (for example, the shared venv at
        \texttt{C:\textbackslash Users\textbackslash Isaac\textbackslash Documents\textbackslash
        TFG\textbackslash Shielded-RL\textbackslash TFG-RL}).
  \item Verify that the Python dependencies are installed (PyTorch with CUDA support,
        Gymnasium, NumPy, Matplotlib, Streamlit; SQLite is part of the Python standard
        library).
\end{enumerate}

\begin{figure}[H]
  \centering
  \includegraphics[width=0.9\textwidth]{imgs/code/repro_launch}
  \caption{CARLA server launch command. The \texttt{-RenderOffScreen} flag disables the GPU
           renderer for headless training; \texttt{-carla-port=2000} sets the TCP port.}
  \label{fig:repro-launch}
\end{figure}

\section{Training commands}\label{app:repro-training}

Figure~\ref{fig:repro-training} shows the commands used to reproduce the three training
configurations of Chapter~\ref{ch:experimentation}.

\begin{figure}[H]
  \centering
  \includegraphics[width=0.9\textwidth]{imgs/code/repro_training}
  \caption{Training commands for the baseline (no shield), basic shield, adaptive shield and
           curriculum-ablation configurations. Each run writes metrics to
           \texttt{runs/<model>\_<shield>\_<timestamp>/metrics.sqlite} and checkpoints to
           \texttt{data/models/<model>\_<shield>\_\{best,final\}.pth}.}
  \label{fig:repro-training}
\end{figure}

\section{Evaluation commands}\label{app:repro-evaluation}

Figure~\ref{fig:repro-evaluation} shows the evaluation commands, both in headless mode (metrics
only) and with the Matplotlib dashboard.

\begin{figure}[H]
  \centering
  \includegraphics[width=0.9\textwidth]{imgs/code/repro_evaluation}
  \caption{Evaluation commands. The shield configuration must match the one used during
           training, as a policy trained with a shield has learned to rely on its interventions.}
  \label{fig:repro-evaluation}
\end{figure}

\section{Live dashboard}\label{app:repro-dashboard}

Figure~\ref{fig:repro-dashboard} shows how to launch the Streamlit dashboard. The dashboard
reads every \texttt{runs/*/metrics.sqlite} found under the project root and allows the user
to select the run(s) to visualise. Because the database is opened in WAL mode, the dashboard
can be launched while a training run is still in progress.

\begin{figure}[H]
  \centering
  \includegraphics[width=0.6\textwidth]{imgs/code/repro_dashboard}
  \caption{Streamlit dashboard launch command.}
  \label{fig:repro-dashboard}
\end{figure}

\section{Query recipes}\label{app:repro-queries}

The following queries document how the empirical numbers in the tables of
Chapter~\ref{ch:experimentation} are computed from the SQLite databases.

\begin{figure}[H]
  \centering
  \includegraphics[width=0.9\textwidth]{imgs/code/repro_sql_rates}
  \caption{SQL queries for the final 100-episode rolling reward and safety rates (crash and
           off-road). Replace the \texttt{key} strings with the metric of interest.}
  \label{fig:repro-sql-rates}
\end{figure}

\begin{figure}[H]
  \centering
  \includegraphics[width=0.9\textwidth]{imgs/code/repro_sql_risklevels}
  \caption{SQL query for the fraction of steps spent in each risk level (adaptive shield only).}
  \label{fig:repro-sql-risklevels}
\end{figure}

\begin{figure}[H]
  \centering
  \includegraphics[width=0.9\textwidth]{imgs/code/repro_sql_interventions}
  \caption{SQL query for the average number of shield interventions per episode by threat
           category (adaptive shield only).}
  \label{fig:repro-sql-interventions}
\end{figure}

\begin{figure}[H]
  \centering
  \includegraphics[width=0.9\textwidth]{imgs/code/repro_sql_curriculum}
  \caption{SQL query for curriculum stage transitions. The \texttt{LAG} window function
           identifies advance and rollback events by comparing consecutive NPC counts.}
  \label{fig:repro-sql-curriculum}
\end{figure}

\section{Checkpoint inspection}\label{app:repro-ckpt}

The utility class \texttt{ModelManager} in \texttt{utils/} exposes helpers to list, compare and
inspect checkpoints. Figure~\ref{fig:repro-modelmanager} shows a quick sanity-check invocation.

\begin{figure}[H]
  \centering
  \includegraphics[width=0.75\textwidth]{imgs/code/repro_modelmanager}
  \caption{Checkpoint inspection with \texttt{ModelManager}. The \texttt{print\_model\_info}
           call reports the network architecture, checkpoint format version, training seed, and
           the state of the running-mean normaliser.}
  \label{fig:repro-modelmanager}
\end{figure}

\section{Unit tests}\label{app:repro-tests}

The reward shaper tests (Figure~\ref{fig:repro-pytest}) run without a CARLA server. They exercise
each component of the shaped reward of Equation~\eqref{eq:shaped-reward} against a crafted
\texttt{info} dictionary and ensure that disabling a weight zeroes out the corresponding
component. They are run locally before every commit that touches \texttt{src/reward\_shaper.py}.

\begin{figure}[H]
  \centering
  \includegraphics[width=0.7\textwidth]{imgs/code/repro_pytest}
  \caption{Running the reward-shaper unit tests. The suite completes in under one second and
           requires no running CARLA server.}
  \label{fig:repro-pytest}
\end{figure}
